\begin{longtable}{ll}
	\textbf{Kod} & \textbf{Fråga / Svar}                                                \\ \hline
	\endhead
	\caption{Q-koder}                                                                   \\
\endlastfoot
	QRA          & Vad heter er station?                                                \\
                     & Vår station heter ...                                                \\ \hline
	QRB          & Hur långt bort från min station befinner ni er?                      \\
                     & Avståndet mellan oss är ungefär ...                                  \\ \hline
	QRG          & Kan ni ange min exakta frekvens?                                     \\
                     & Er exakta frekvens är ... (MHz/kHz)                                  \\ \hline
	QRH          & Varierar min frekvens/våglängd?                                      \\
                     & Er frekvens/våglängd varierar.                                       \\ \hline
	QRI          & Hur är min sändningston (CW)?                                        \\
                     & Er sändningston är 1--God, 2--Varierande, 3--Dålig                   \\ \hline
	QRK          & Vilken uppfattbarhet har mina signaler?                              \\
                     & Uppfattbarheten hos dina signaler är:                                \\
                     & 1--Dålig, 2--Bristfällig, 3--Ganska god, 4--God, 5--Utmärkt          \\ \hline
	QRL          & Är ni upptagen?                                                      \\
                     & Jag är upptagen med ... (namn/signal) stör ej.                       \\ \hline
	QRM          & Är ni störd av annan station?                                        \\
                     & Störningarna är:                                                     \\
                     & 1--Obef., 2--Svaga, 3--Måttliga, 4--Starka, 5--Mycket starka         \\ \hline
	QRN          & Besväras ni av atmosfäriska störningar?                              \\
                     & Störningarna är:                                                     \\
                     & 1--Obef., 2--Svaga, 3--Måttliga, 4--Starka, 5--Mycket starka         \\ \hline
	QRO          & Kan jag (ska jag) öka sändareffekten?                                \\
                     & Öka sändareffekten.                                                  \\ \hline
	QRP          & Kan jag (ska jag) minska sändareffekten?                             \\
                     & Minska sändareffekten.                                               \\ \hline
	QRQ          & Kan jag (får jag) öka sändningshastigheten?                          \\
                     & Öka sändningshastigheten.                                            \\ \hline
	QRS          & Kan jag (skall jag) sända långsammare?                               \\
                     & Sänd långsammare.                                                    \\ \hline
	QRT          & Skall jag avbryta sändningen?                                        \\
                     & Avbryt sändningen                                                    \\ \hline
	QRU          & Har ni något till mig?                                               \\
                     & Jag har inget till er. Se även QTC.                                  \\ \hline
	QRV          & Är ni redo?                                                          \\
                     & Jag är redo.                                                         \\ \hline
	QRX          & När anropar ni mig härnäst?                                          \\
                     & Jag anropar er kl ... (på ... MHz/kHz)                               \\ \hline
	QRZ          & Vem anropar mig?                                                     \\
                     & Ni anropas av ... (på ... MHz/kHz).                                  \\ \hline
	QSA          & Vilken styrka har mina signaler?                                     \\
                     & Era signaler är:                                                     \\
                     & 1--Ej uppf., 2--Svaga, 3--Ganska starka, 4--Starka, 5--Mycket starka \\ \hline
	QSB          & Svajar styrkan på mina signaler?                                     \\
                     & Styrkan på era signaler svajar.                                      \\ \hline
	QSK          & Kan du höra mig mellan dina tecken och får jag avbryta dig?          \\
                     & Jag kan höra dig mellan mina tecken och du får avbryta.              \\ \hline
	QSL          & Kan ni ge mig kvittens?                                              \\
                     & Jag kvitterar.                                                       \\ \hline
	QSO          & Ha ni förbindelse med ... eller ... (förmedlat)?                     \\
                     & Jag har förbindelse med ... (via ...)                                \\ \hline
	QST          & Har tidigare använts som allmänt anrop men ersatts av CQ             \\ \hline
	QSY          & Skall jag övergå till att sända på annan frekvens?                   \\
                     & Gå över till att sända på annan frekvens (eller ... kHz/MHz).        \\ \hline
	QTC          & Hur många telegram har ni att sända?                                 \\
                     & Jag har ... telegram till dig (eller ...).                           \\ \hline
	QTH          & Vilken är er geografiska position?                                   \\
                     & Min geografiska position är ...                                      \\ \hline
	QTR          & Kan ni ge mig rätt tid?                                              \\
                     & Rätt tid är ...                                                      \\
\end{longtable}